% =====================================================================
%  ch5_equivalence.tex — §5 等价构造
% =====================================================================
\section{等价的构造}\label{s:equivalence}

\subsection{给 $\cF(M)$ 的对象配置向量丛}\label{ss:assignbundles}

回忆 $\cF(M)$ 中 v.b. 型对象是一对 $(L,\cE)$：$L$ 是以 $V(G)$ 中顶点为顶点、在每个 $M_e$ 中都有边的图，$\cE$ 是 $L$ 上自由有限生成 $\cO_K$-模的局部系统。我们固定这样一个拉格朗日图 $L_0$，使得对 $(L_0,\cO_K)$ 将对应结构层 $\cO_X$。

现在对任意对象 $(L,\cE)\in\cF(M)$ 与每条连接 $v,v'\in V(G)$ 的边 $e\in E(G)$，有下列不变量：
\begin{itemize}
\item $r_e(L)=r_e(L_0,L)\in \Z$，即 $M_e$ 中 $L$ 关于 $L_0$ 沿负方向的旋转数。和 $\sum_{e\in E(G)}r_e(L)$ 将是相应向量丛的斜率。
\item $S_{e/v}(L)=S_{e/v}(L_0,L)$，即在 $M_e\setminus\{p_v,p_{v'}\}$ 的万有覆盖上由 $L_0$ 与 $L$ 围出的带符号面积，其中取在接近 $p_{v'}$ 时彼此靠近的提升。有 $S_{e/v'}(L)+S_{e/v}(L)=r_{e}(L)A_e$。
\item 单值性 $R_{\cE,e/v}:\cE_v\to\cE_{v'}$。
\end{itemize}

如下定义 $X$ 上的向量丛 $\Phi(L,\cE)$。首先，它拉回到 $Y_v-F_v$ 上由 $\cE_v\otimes_{\cO_K}\cO_{Y_v-F_v}$ 给出。接着要描述群胚 $\pi_1(G)$ 的作用。对每条看作 $\pi_1(G)$ 中态射 $v\to v'$ 的边 $e$，相应的同构
\[
u_{e/v}:\cE_v\otimes\cO_{Y_v-F_v}\to g_{e/v}^*(\cE_{v'}\otimes\cO_{Y_{v'}-F_{v'}})
\]
由
\begin{equation}\label{eq:gluingdata}
u_{e/v}=R_{\cE,e/v}\otimes T^{-S_{e/v}(L)}t_{e/v}^{r_e(L)}
\end{equation}
给出。若 A 模型作了体形变，此式应乘以 $\mathfrak{b}$ 在 $M_e\setminus\{p_v,p_{v'}\}$ 的万有覆盖上由 $L_0$ 与 $L$ 围出的面积上的积分之指数。

由 \cite{Fa}，向量丛 $\Phi(L,\cE)$ 半稳定，斜率为 $\sum_{e\in E(G)}r_e(L)$。

\subsection{复形的抽象同调扰动引理（HPL）}\label{ss:HPL}

回忆下面同调扰动的抽象框架，细节见 \cite{CL}。为简单起见，基域取 Novikov 域 $K$。

设 $(\cK,d_{\cK})$ 与 $(\cL,d_{\cL})$ 是复形。设有映射 $i,p,h$：$i:\cL\to\cK$ 与 $p:\cK\to \cL$ 是复形态射，$h:\cK\to\cK$ 是（上同调）次数 $-1$ 的映射，满足
\[
pi=1_{\cL},\qquad 1_{\cK}-ip=dh+hd,\qquad h^2=0,\qquad ph=0,\qquad hi=0.
\]

再取微分 $d_{\cK}$ 的扰动 $\delta$，满足 Maurer--Cartan 方程 $[d_{\cK},\delta]+\delta^2=0$。于是 $\tilde{d}_{\cK}=d_{\cK}+\delta$ 是 $\cK$ 上的微分。假定 $\cK$ 的自同态 $(\id_{\cK}+h\delta)$ 可逆（从而 $(\id_{\cK}+\delta h)$ 也可逆）。那么 $d_{\cL}$、$i$、$p$、$h$ 都有自然的扰动，使上述所有关系继续成立：
\[
\tilde{d}_{\cL}=d_{\cL}+p\delta(\id+h\delta)^{-1}i,\quad
\tilde{i}=(\id+h\delta)^{-1}i,\quad
\tilde{p}=p(\id+\delta h)^{-1},\quad
\tilde{h}=(\id+h\delta)^{-1}h.
\]
特别地，$\tilde{i}$ 与 $\tilde{p}$ 是带扰动微分的复形之间的拟同构。

\begin{remark}\label{remark:top_nilpotent}
假定 $\cK$ 与 $\cL$ 的（分次）分支是 $K$ 上的 Banach 向量空间，且 $i,p,h,\delta$ 连续。此时 $(\id+h\delta)$ 可逆这一假设可由 $h\delta:\cK\to\cK$ 局部拓扑幂零推出，即对任意齐次 $x\in\cK$ 有 $\lim_{n\to\infty}(h\delta)^n(x)=0$。事实上此时
\[
(\id+h\delta)^{-1}(x)=\sum_{n=0}^{\infty}(-1)^n(h\delta)^n(x),\qquad
(\id+\delta h)^{-1}(x)=\sum_{n=0}^{\infty}(-1)^n(\delta h)^n(x).
\]
于是 $\tilde{d}_{\cL},\tilde{i},\tilde{p},\tilde{h}$ 的公式也可写成无穷和。
\end{remark}

注意，对如上两组数据 $(\cK,\cL,i,p,h)$ 与 $(\cK',\cL',i',p',h')$，可定义其张量积为 $(\cK\otimes\cK',\cL\otimes\cL',i\otimes i', p\otimes p',h'')$，其中 $h''=h\otimes\id+ip\otimes h'$。

现在设 $(\cK,\cL,i,p,h)$ 如上，且 $\mu_{\cK}=(\mu_{\cK}^1,\mu_{\cK}^2,\dots)$ 是 $\cK$ 上的 $A_{\infty}$-结构。利用张量积公式便得到如上数据 $(T(\cK[1]),T(\cL[1]),i',p',h')$。我们得到次数 $1$ 的余导子 $\delta:T(\cK[1])\to T(\cK[1])$，其分量为 $\delta^1=\mu^1_{\cK}-d_{\cK},\delta^2=\mu^2_{\cK},\dots$。假定 $(\id_{\cK}+h\delta^1)$ 可逆，则易见 $(\id_{T(\cK[1])}+h'\delta)$ 也可逆。应用上述公式得到变形微分 $\tilde{d}_{T(\cL[1])}$，它实际上是余导子，因此给出 $\cL$ 上的 $A_{\infty}$-结构 $\mu_{\cL}$。变形后的态射 $\tilde{i'}:T(\cL[1])\to T(\cK[1])$、$\tilde{p'}:T(\cK[1])\to T(\cL[1])$ 实际上是 DG 余代数的态射，因此给出 $A_{\infty}$-代数态射 $\tilde{i}:(\cL,\mu_{\cL})\to (\cK,\mu_{\cK})$、$\tilde{p}:(\cK,\mu_{\cK})\to (\cL,\mu_{\cL})$，它们都是拟同构。细节见 \cite[Section 3.3]{CL}。

注意若处于注~\ref{remark:top_nilpotent} 的情形，则 $\mu_{\cL},\tilde{i},\tilde{p}$ 作为无穷和的表达式正是标准的树和。若 $\delta^1=0$，则分量 $\mu_{\cL}^n,\tilde{i}_n,\tilde{p}_n$ 的求和是有限的。

同样的构造也适用于 $A_{\infty}$-范畴。我们将在第~\ref{ss:HPLFloer} 节用它论证：涉及「无穷」Hamilton 扰动的 A 模型变体 $\cF(M;H)$ 与 $\cF(M)$ 拟同构；也用它得到对应于 Floer 复形生成元的 theta 函数表达式（只要这些复形集中在次数零）。

\subsection{无穷 Hamilton 扰动与 \v{C}ech 复形}\label{ss:cech}

回忆命题~\ref{prop:ainfty_ideal}：记 $P_v$ 为各半边 $e/v$ 对应的子集 $([0,3]\times S^1/\!\sim)\subset M_e$ 之并，$N_e$ 为子集 $[1,3]\times S^1\subset M_e$。则位于 $P_v$（相应地 $N_e$）之外的 Floer 复形生成元（在完备化后）张成 $\cF(M;H)$ 中的 $A_\infty$-理想，记 $\cF(P_v;H)$（相应地 $\cF(N_e;H)$）为 $\cF(M;H)$ 对此 $A_\infty$-理想的商范畴。

这些商带有 $A_\infty$-函子 $\cF(M;H)\to \cF(P_v;H)$ 与 $\cF(P_v;H)\to \cF(N_e;H)$，它们在态射上满射、高阶项为零。因此朴素的链层极限全忠实地嵌入同伦极限
\[
\lim\Bigl(\prod_{v\in V(G)}\cF(P_v;H)\prarr\prod_{e\in E(G)}\cF(N_e;H)\Bigr)
\hto
\holim\Bigl(\prod_{v\in V(G)}\cF(P_v;H)\prarr\prod_{e\in E(G)}\cF(N_e;H)\Bigr),
\]
而 Fukaya 范畴 $\cF(M;H)$ 全忠实地嵌入朴素极限
\[
\cF(M;H)\hto \lim\Bigl(\prod_{v\in V(G)}\cF(P_v;H)\prarr\prod_{e\in E(G)}\cF(N_e;H)\Bigr).
\]
我们将在第~\ref{ss:localfunctors} 节证明有自然等价
\[
\Perf(\cF(P_v;H))\simeq \Perf(U_v),\qquad \Perf(\cF(N_e;H))\simeq \Perf(U_e),
\]
在这些等价下函子 $\cF(P_v;H)\to \cF(N_e;H)$ 对应限制函子 $\Perf(U_v)\to \Perf(U_e)$。于是得到全忠实嵌入
\[
\cF(M;H)\to\holim\Bigl(\prod_{v\in V(G)}\Perf(U_v)\prarr\prod_{e\in E(G)}\Perf(U_e)\Bigr)\simeq \Perf(X_K),
\]
它诱导全忠实函子 $\Psi:\Perf(\cF(M;H))\to \Perf(X_K)$。但在对象同构类层面上，此函子恰把 $(L,\cE)$ 送到 $\Phi(L,\cE)$。由于形如 $\Phi(L,\cE)$ 的向量丛生成范畴 $\Perf(X_K)$，故 $\Psi$ 是等价。

\subsection{对 Hamilton 扰动的 HPL}\label{ss:HPLFloer}

现在说明 HPL 如何给出 $A_{\infty}$-范畴 $\cF(M)$ 与 $\cF(M;H)$ 之间的拟等价，并解释它如何看作第~\ref{ss:continuation} 节所述 continuation 函子 $\mathfrak{K}:\cF(M)\to \cF(M;H)$ 的代数对应物。

与第~\ref{ss:FMH} 节一样，对连接顶点 $v$ 与 $v'$ 的每条边 $e$，把 $M_e$ 与 $[0,4]\times(\R/\Z)/\!\sim$ 等同，坐标为 $(\tau,\psi)$（$\tau=0$ 对应节点 $p_v$，$\tau=4$ 对应 $p_{v'}$）。考虑两个 v.b. 型拉格朗日 $L,L'$。不妨设 $CF^*(L,L';\varepsilon h)$ 的生成元远离扰动 $f_n$ 与 $f$ 的支集，且 $\cX(L,L';H_n)$（相应地 $\cX(L,L';H)$）与 $\cX(L,L';\varepsilon h)$ 的差别只是每个分量 $M_e$ 中多出下列生成元：
\begin{itemize}
\item $n$ 个（相应地无穷多个）次数 $1$ 的生成元 $q_{e/v,1},q_{e/v,2},\dots$（按 $\tau$ 坐标递增排列），位于 $(0,1)\times S^1$ 中；
\item $2n$ 个（相应地「两倍」无穷多个）次数 $0$ 的生成元 $\dots,p_{e/v,2},p_{e/v,1}$（靠近 $\tau=1$）与 $p_{e/v',1},p_{e/v',2},\dots$（靠近 $\tau=3$），位于 $(1,3)\times S^1$ 中；
\item $n$ 个（相应地无穷多个）次数 $1$ 的生成元 $\dots,q_{e/v',2},q_{e/v',1}$，位于 $(3,4)\times S^1$ 中。
\end{itemize}

记 $\mu^1_{nv}$ 为 $CF^*(L,L';H)$ 上的微分 $\mu^1_H$ 的「朴素」（或「低面积」）部分，它只涉及支撑在 $\tau=1$ 或 $\tau=3$ 附近、位于单个分量 $M_e$ 中（无传播）的全纯圆盘。于是 $\mu^1_{nv}$ 把 $p_{e/v,k}$ 送到 $q_{e/v,k}$ 的倍数（对每条半边 $e/v$ 与所有 $k\geq 1$），把其余生成元送到零。连接 $p_{e/v,k}$ 与 $q_{e/v,k}$ 的圆盘面积可通过缩小扰动 $f_n$ 与 $f$ 的支集变得任意小；这使我们能假定 Floer 微分 $\mu^1_H$ 的所有其他贡献面积都大于 $\mu^1_{nv}$ 所记录的那些。

令 $\delta^1=\mu^1_H-\mu^1_{nv}$（并对 $k\geq 2$ 令 $\delta^k=\mu^k_H$），现处于抽象 HPL 的框架中。自然包含 $i:(CF^*(L,L'),0)\to (CF^*(L,L';H),\mu^1_{nv})$ 是复形态射，投影 $p:(CF^*(L,L';H),\mu^1_{nv})\to (CF^*(L,L'),0)$ 也是。再取同伦 $h$ 为把每个新的次数 $1$ 生成元 $q_{e/v,k}$ 送到相应次数 $0$ 生成元 $p_{e/v,k}$ 的映射，乘以 $\mu^1_{nv}$ 中出现的系数之逆。于是 $h\delta^1$ 局部拓扑幂零，因为对 $\delta^1$ 有贡献的扰动全纯圆盘辛面积大于对 $\mu^1_{nv}$ 有贡献的圆盘。由此 $\id+h\delta$ 可逆（见注~\ref{remark:top_nilpotent}），可以应用 HPL。

% [Tutor Note: 点拨] 这一段的几何内容是：$\delta^1=\mu^1_H-\mu^1_{nv}$ 只记录「被扰动拉宽、因而必须穿过 $\tau=1$ 或 $\tau=3$ 处新生成元」的带，$h$ 则把终端 $q_{e/v,k}$ 拉回起点 $p_{e/v,k}$。于是 $(h\delta^1)^\ell$ 逐项对应「连续穿过 $\ell$ 个节点」的传播带，而按面积严格递增，故级数在完备化中收敛；这与 §3.2 中「传播圆盘必须允许穿过节点」的几何观察是同一件事的代数化。依据：原文 §5.4 与 §3.2，并见本工程 01_通读笔记。

把此构造应用于 $A_\infty$-范畴 $\cF(M)$ 与 $\cF(M;H)$（更确切地说，应用于其对象满足上述关于 $\tau=1$、$\tau=3$ 附近无 Floer 生成元、以及 Floer 复形在 Hamilton 扰动下行为的假设的全子范畴），便得到 Floer 复形 $CF^*(L,L')$ 上的运算 $\mu^k_{HPL}$（$k\geq 1$），它们由第~\ref{ss:HPL} 节的公式给出；还得到 $A_\infty$-函子 $\tilde{i}$ 与 $\tilde{p}$，它们给出此 $A_\infty$-范畴与 $\cF(M;H)$ 之间的拟等价。

下面说明由同调扰动理论从 $\mu^k_H$ 得到的运算 $\mu^k_{HPL}$ 等于 Fukaya 范畴 $\cF(M)$ 的结构映射 $\mu^k$，从而 $\tilde{i}$ 与 $\tilde{p}$ 实际上给出 $\cF(M)$ 与 $\cF(M;H)$ 之间的拟等价。

从微分开始。回忆 HPL 给出
\begin{equation}\label{eq:mu1_HPL}
\mu^1_{HPL}=p\delta^1(\id+h\delta^1)^{-1}i=\sum_{\ell=0}^\infty
(-1)^\ell p\delta^1(h\delta^1)^\ell i.
\end{equation}

考虑如上的两个 v.b. 型拉格朗日 $L,L'$，以及一个对 $CF^*(L,L')$ 上 Floer 微分有贡献的传播全纯带 $u:S\to M$，它把输入生成元 $p_1$ 经一串依次位于分量 $M_{e_1},\dots,M_{e_\ell}$ 中的全纯带（其中 $p_1\in M_{e_1}$、$p_0\in M_{e_\ell}$）连到输出生成元 $p_0$，这些带通过节点 $p_{v_1},\dots,p_{v_{\ell-1}}$ 彼此粘接。由于 $u$ 刚性，其边界沿 $L$ 与 $L'$ 行进而不折返，且 $\tau$ 坐标在每个分量上单调变化。把每条边 $e_j$ 定向，使带沿 $M_{e_j}$ 从输入到输出按 $\tau$ 增大方向行进，即 $p_{v_j}$ 位于 $M_{e_j}$ 的 $\tau=4$ 端与 $M_{e_{j+1}}$ 的 $\tau=0$ 端。例如设输入 $p_1\in M_{e_1}$ 的 $\tau$ 坐标小于 $1$，输出 $p_0\in M_{e_\ell}$ 的 $\tau$ 坐标大于 $1$，于是 $u$ 的每个分量都穿过圆周 $\{1\}\times S^1\subset M_{e_i}$（其余情形类似）。记 $w_j\in \R_+$ 为 $u$ 的第 $j$ 个分量在 $\tau=1$ 处的宽度，即带在 $M_{e_j}-\{p_{v_{j-1}},p_{v_j}\}$ 的万有覆盖上提升后，两条边界在 $\tau=1$ 处 $\psi$ 坐标之差；令 $k_j=\lceil w_j \rceil$。于是 Hamilton 扰动 $H$（或对 $n>\max(w_j)$ 的 $H_n$）把 $u$ 的每个分量裂成一个在 $\tau$ 到达 $1$ 之前终止于新的次数 $1$ 生成元 $q_{e_j/v_{j-1},k_j}$ 的带，以及一个刚从 $\tau=1$ 之后由新的次数 $0$ 生成元 $p_{e_j/v_{j-1},k_j}$ 开始的带。这样便可把 $u$ 联系到一串 $\ell+1$ 个对 $CF^*(L,L';H)$ 上微分 $\mu^1_H$（从而对 $\delta^1=\mu^1_H-\mu^1_{nv}$）有贡献的扰动传播全纯带，其间夹着 $\ell$ 条生成元对 $p_{e_j/v_{j-1},k_j}$ 与 $q_{e_j/v_{j-1},k_j}$ 之间的低面积连接轨道。这些正是式~\eqref{eq:mu1_HPL} 右端所计数的构形。此外，$u$ 的传播重数等于它裂成的 $\ell+1$ 个扰动带传播重数之积；其面积等于这些带面积之和减去生成元对 $p_{e_j/v_{j-1},k_j}$ 与 $q_{e_j/v_{j-1},k_j}$ 之间连接轨道面积之和，和乐也类似。最后，符号 $(-1)^\ell$ 来自破裂构形中在新的次数 $1$ 生成元 $q_{e_j/v_{j-1},k_j}$ 处多出的输出对的整体符号贡献：每次终止于 $q_{e_j/v_{j-1},k_j}$ 的两条轨道沿 $L'$ 的边界定向相反，故符号相差 $-1$。由此 $\mu^1_{HPL}=\mu^1$。

对 $k\geq 2$ 的 $\mu^k$ 论证类似。考虑两两满足上述关于扰动下 Floer 复形行为之简化假设的 v.b. 型拉格朗日 $L_0,\dots,L_k$，以及对 $\mu^k$ 有贡献、边界在 $L_0,\dots,L_k$ 上的刚性传播全纯圆盘 $u:S\to M$。$u$ 与 $M$ 各分量中圆周 $\{1\}\times S^1$ 与 $\{3\}\times S^1$ 的邻域 $\cN_\delta$ 之交是传播圆盘的一些带状部分的并。其中沿 $\tau$ 减小（相应地增大）方向穿过 $\tau=1$（相应地 $\tau=3$）的带基本不受 Hamilton 扰动 $H_n$ 影响，而沿 $\tau$ 增大（相应地减小）方向穿过 $\tau=1$（相应地 $\tau=3$）的带，一旦 $n$ 超过其沿 $\psi$ 坐标的宽度，就按上述方式裂开。于是 $u:S\to M$ 裂成若干对 $\cF(M;H)$ 结构映射有贡献的扰动传播圆盘之并，每个的输入或是 $u$ 的输入（故是来自 $\cX(L_{i-1},L_i;\varepsilon h)$ 的「旧」生成元），或是新的次数 $0$ 生成元 $p_{e/v,k}$；输出或是新的次数 $1$ 生成元 $q_{e/v,k}$，或是 $u$ 的输出。因为 $h$ 除新的次数 $1$ 生成元外处处为零，并把后者送到相应新的次数 $0$ 生成元，这类构形恰与 HPL 公式中出现的树和一致，于是 $\mu^k_{HPL}=\mu^k$。

这就完成了 $\cF(M;H)$ 与 $\cF(M)$ 拟等价的证明（经由 HPL 给出的 $A_\infty$-函子 $\tilde{i}$ 与 $\tilde{p}$）。

虽然论证中并不需要，比较 $\tilde{i}:\cF(M)\to \cF(M;H)$ 与第~\ref{ss:continuation} 节所述 continuation 函子 $\mathfrak{K}$ 仍有启发。线性项的 HPL 公式为
\[
\tilde{i}^1=(\id+h\delta^1)^{-1}i=\sum_{\ell=0}^\infty(-1)^\ell (h\delta^1)^\ell i.
\]
由于 $h\delta^1$ 在次数 $1$ 生成元上为零，对 $CF^1$ 此式化简为朴素包含 $i$。对次数 $0$ 的生成元，$\tilde{i}^1$ 与包含之差是一些破裂构形的计数，这些构形由对微分 $\mu^1_H$（即对 $\delta^1=\mu^1_H-\mu^1_{nv}$）有贡献、终止于次数 $1$ 生成元 $q_{e_j/v_{j-1},k}$ 的扰动传播全纯带组成，其间夹着生成元对 $q_{e_j/v_{j-1},k}$ 与 $p_{e_j/v_{j-1},k}$ 之间低面积连接轨道之逆。如前论证，这类构形几乎恰好对应 Floer 微分 $\mu^1$（Hamilton 扰动为 $\varepsilon h$）的传播全纯圆盘，只有承载输出的分量不同；后者可直接计算，验证其与从 Hamilton 扰动 $\varepsilon h$ 到扰动 $H_n$（$n$ 充分大）的 Floer continuation 轨道行为一致。与注~\ref{rmk:continuation} 的描述比较，得 $\tilde{i}^1=\mathfrak{K}^1$。这又表明 $\mathfrak{K}$ 是拟等价。我们预期（但未验证）$A_\infty$-函子 $\tilde{i}$ 与 $\mathfrak{K}$ 的高阶项也一致。

\subsection{局部函子}\label{ss:localfunctors}

现在描述函子 $\cF(P_v;H)\to \Perf(U_v)$ 与 $\cF(N_e;H)\to \Perf(U_e)$，它们粘起来便给出函子 $\cF(M;H)\to\Perf(X_K)$。

从 $P_v$ 开始。我们把每个 v.b. 型对象 $(L,\cE)$ 送到自由层 $\cE_{v}\otimes_{\cO_K}\cO_{U_v}$。

设 $(L,\cE),(L',\cE')$ 是两个 v.b. 型对象。由于局部系统 $\cE$、$\cE'$ 在 $P_v$ 上可平凡化，记号和论证中略去它们，假定处理的是平凡秩 $1$ 局部系统。又暂设位于 $P_v$ 内的 $\cX(L,L';\varepsilon h)$ 的唯一元素就是节点 $p_v$ 本身（这总可通过 Hamilton 同痕达到），且 $P_v$ 内 $\cX(L,L';H)$ 的元素由生成元 $p_v$ 以及（对每条半边 $e/v$）无穷多个次数 $1$ 的生成元 $q_{e/v,0},q_{e/v,1},q_{e/v,2},\dots$ 与次数 $0$ 的生成元 $\dots,p_{e/v,2},p_{e/v,1},p_{e/v,0},p_{e/v,-1},p_{e/v,-2},\dots$ 组成，它们位于 $(0,3)\times S^1\subset M_e$ 中。（次数 $0$ 生成元的编号使得 $p_{e/v,k}$ 在 $k\geq 0$ 时靠近 $\tau=1$，在 $k<0$ 时靠近 $\tau=3$。）

Floer 微分把 $p_v$ 送到沿三条边紧邻的三个次数 $1$ 生成元 $q_{e/v,0}$ 的线性组合，把每个 $p_{e/v,k}$（$k\geq 0$）送到相应生成元 $q_{e/v,k}$ 的倍数（这些不涉及传播）。它还把 $p_{e/v,-k}$（$k\geq 1$）送到
\[
\sum_{e'/v,\ e'\neq e}\ \sum_{\ell\geq 0} C^{v;e,e'}_{k,\ell}\,
T^{S_{e/v}(p_{e/v,-k})-S_{e'/v}(q_{e'/v,\ell})}\,q_{e'/v,\ell},
\]
其中求和项对应自 $M_e$ 经 $p_v$ 传播到 $M_{e'}$、输入次数为 $k$、输出次数为 $\ell$ 的带；这里 $C^{v;e,e'}_{k,\ell}$ 如定义~\ref{defi:propcoeff}，而 $S_{e/v}(p_{e/v,-k})$ 与 $-S_{e'/v}(q_{e'/v,\ell})$ 是两部分的面积。（按约定，$S_{e/v}(x)$ 记 $M_e$ 内 Floer 生成元 $x$ 到 $p_v$ 的圆盘的带符号面积，故从 $p_v$ 到 $x$ 的圆盘带符号面积为 $-S_{e/v}(x)$。）

由此可知 Floer 微分满射（即使完备化后也如此，因为 Hamilton 函数 $H_n$ 与 $H$ 的构造保证连接 $p_{e/v,k}$ 与 $q_{e/v,k}$ 的轨道面积不依赖 $k$ 而一致有界），且上同调集中在次数零，生成元为
\begin{eqnarray}\label{eq:tildepv}
\tilde{p}_v&=&p_v\,+\,\sum_{e/v}\,
T^{-S_{e/v}(p_{e/v,0})}\,p_{e/v,0}\quad \text{与}\\
\label{eq:tildepev}
\tilde{p}_{e/v,-k}&=&p_{e/v,-k}\,+\sum_{e'/v,\ e'\neq e}\ \sum_{\ell\geq 0} C^{v;e,e'}_{k,\ell}\,
T^{S_{e/v}(p_{e/v,-k})-S_{e'/v}(p_{e'/v,\ell})}\,p_{e'/v,\ell},
\end{eqnarray}
其中 $T$ 的指数对应 $p_v$ 与相应生成元之间轨道的面积。对一般的 v.b. 型对象，适当重新编号生成元后情形类似。

事实上有一个简单的几何模型，记作 $\cF(P_v)$，其中 Floer 微分消失，态射空间是 $\cF(P_v;H)$ 中态射空间的上同调。具体说，取一个在 $P_v$ 内部像 $\varepsilon h$、在 $P_v$ 边界附近（即交于 $p_v$ 的三个 $M$ 分量中各 $\tau=3$ 处）像 $H$ 的 Hamilton 函数。由与第~\ref{ss:FMH} 节相同的论证，$P_v$ 外的生成元仍构成 $A_\infty$-理想；而 $P_v$ 内的生成元现由 $p_v$ 与所有 $p_{e/v,-k}$（$e/v$ 任意、$k\geq 1$）组成，都在次数零。无论用 HPL 还是用 continuation 映射，都可见 $\cF(P_v)$ 与 $\cF(P_v;H)$ 拟等价，且拟等价的线性项把 $p_v$ 送到 $\tilde{p}_v$，把 $p_{e/v,-k}$ 送到 $\tilde{p}_{e/v,-k}$。

现在如何定义从 $\cF(P_v)$（相应地 $\cF(P_v;H)$）到 $\Perf(U_v)$ 的、作用在 v.b. 型对象间态射空间（相应地闭的次数零态射）上的函子已很清楚：把 $p_v$（相应地 $\tilde{p}_v$）送到 $U_v$ 上的常函数 $1$，把 $p_{e/v,-k}$（相应地 $\tilde{p}_{e/v,-k}$）送到
\[
T^{S_{e/v}(p_{e/v,-k})}\, t_{e/v}^{-k}.
\]
为证这确是函子，需验证：$\cF(P_v)$ 中的 Floer 乘积，对应 $U_v$ 上函数的乘积。记 $\hat{p}_{e/v,-k}=T^{-S_{e/v}(p_{e/v,-k})}\,p_{e/v,-k}$ 为按面积权重新标度后的 Floer 生成元。考察 $P_v$ 中的传播全纯圆盘，其两个输入取在不同分量（$e_1\neq e_2$）上的生成元 $p_{e_1/v,-k_1}$ 与 $p_{e_2/v,-k_2}$，则
\[
\mu^2(\hat{p}_{e_1/v,-k_1},\hat{p}_{e_2/v,-k_2})=K^{v;e_1,e_2}_{k_1,k_2} p_v+
\sum_{b=0}^{k_1-1}C^{v;e_2,e_1}_{k_2,b} \hat{p}_{e_1/v,b-k_1} +
\sum_{a=0}^{k_2-1}C^{v;e_1,e_2}_{k_1,a} \hat{p}_{e_2/v,a-k_2},
\]
这与式~\eqref{eq:partialproduct} 的乘积公式完全一致。而对同一分量上的生成元，由于 $\mu^2(\hat{p}_{e/v,-k_1},\hat{p}_{e/v,-k_2})=\hat{p}_{e/v,-k_1-k_2}$，结论立得。

（若要在 $\cF(P_v;H)$ 的态射空间其余部分显式定义函子，最好用同调扰动理论把严格函子 $\cF(P_v)\to \Perf(U_v)$ 提升为 $A_\infty$-函子 $\cF(P_v;H)\to \Perf(U_v)$；不过我们不需要显式公式。）

最后，验证函子全忠实涉及完备化的比较。$\cF(P_v)$ 中的态射是 Floer 生成元的无穷线性组合，其系数的 Novikov 赋值趋于 $+\infty$；而开仿射域 $U_v$ 上的函数是基函数 $1$ 与 $t_{e/v}^{-k}$（$e/v$ 任意、$k\geq 1$）的线性组合，收敛性在 $|t_{e/v}|\geq |q_e|^{3/4}$（即 $\val(t_{e/v})\leq \frac34 A_e$）时成立。在我们的函子把 $p_{e/v,-k}$ 送到 $T^{S_{e/v}(p_{e/v,-k})}\,t_{e/v}^{-k}$ 之下，两个完备化一致，这直接来自下述几何事实：连接 $M_e$ 中 $\tau=3$ 附近生成元 $p_{e/v,-k}$ 与 $p_v$ 的次数 $k$ 圆盘面积 $S_{e/v}(p_{e/v,-k})$ 与 $\frac34 k A_e$ 相差有界。

函子 $\cF(N_e;H)\to \Perf(U_e)$ 类似构造：所有 v.b. 型对象送到 $U_e$ 上的自由层，Floer 生成元送到坐标 $t_{e/v}$（或另一顶点处的 $t_{e/v'}$）的适当幂次乘以合适因子。把 $N_e$ 看作 $P_v$ 的子集，并考虑一对如前所述在 $P_v$ 中除节点 $p_v$ 外不相交的 v.b. 型拉格朗日，它们在 $\cF(N_e;H)$ 中的态射空间是生成元无穷序列 $p_{e/v,k}$（$k\in\Z$，都在次数零）张成空间的完备化；我们把每个 $p_{e/v,k}$ 送到 $T^{S_{e/v}(p_{e/v,k})}\,t_{e/v}^k$。在此函子下完备化一致，仍来自观察：$S_{e/v}(p_{e/v,k})$ 在 $k\ll 0$（靠近 $\tau=3$ 的生成元）时接近 $\frac34 |k| A_e$，在 $k\gg 0$（靠近 $\tau=1$ 的生成元）时接近 $-\frac14 k A_e$。

按定义，限制函子 $\cF(P_v;H)\to \cF(N_e;H)$ 把态射空间彼此相送的方式就是商去所有位于 $N_e$ 之外的生成元；记此商映射为 $Q$。把它与 HPL（或 continuation）提供的从 $\cF(P_v)$ 到 $\cF(P_v;H)$ 的拟等价复合，得到限制函子 $\cF(P_v)\to \cF(N_e;H)$。由式~\eqref{eq:tildepv}--\eqref{eq:tildepev}，它把 $p_v$ 送到 $Q(\tilde{p}_v)=T^{-S_{e/v}(p_{e/v},0)} p_{e/v,0}$，把 $p_{e/v,-k}$ 送到 $Q(\tilde{p}_{e/v,-k})=p_{e/v,-k}$ 本身，而对 $e'\neq e$，把 $p_{e'/v,-k}$ 送到
\[
Q(\tilde{p}_{e'/v,-k})=\sum_{\ell\geq 0} C^{v;e',e}_{k,\ell}\,
T^{S_{e'/v}(p_{e'/v,-k})-S_{e/v}(p_{e/v,\ell})}\,p_{e/v,\ell}.
\]
利用 $t_{e'/v}^{-k}=\sum_{\ell=0}^\infty C_{k,\ell}^{v;e',e}\, t_{e/v}^\ell$，容易验证这些公式与从 $\Perf(U_v)$ 到 $\Perf(U_e)$ 的限制一致。
